\section{Contexto}

Los mercados de criptoactivos y las finanzas descentralizadas han pasado de ser aplicaciones aisladas sobre blockchain a formar ecosistemas digitales interdependientes. En ellos conviven exchanges centralizados (CEX), exchanges descentralizados (DEX), Automated Market Makers (AMM), protocolos de préstamo, derivados perpetuos, bridges, oráculos, validadores, liquidadores, proveedores de liquidez y agentes automatizados. Bitcoin demostró la viabilidad de un libro mayor distribuido sin una autoridad central \cite{nakamoto2008bitcoin}; Ethereum extendió esa idea con contratos inteligentes programables \cite{buterin2014ethereum,wood2014ethereum}. Esa programabilidad permitió que el intercambio, el préstamo y la provisión de liquidez fueran componentes ejecutables de una infraestructura financiera abierta \cite{schar2021defi,werner2021sok}.

El interés de esta tesis no es recomendar estrategias de inversión ni construir un bot orientado a rentabilidad. El problema de investigación se sitúa en la ingeniería de sistemas: comprender cómo la interacción entre infraestructura computacional, reglas de protocolo, liquidez, latencia e incentivos económicos produce fenómenos que no pueden explicarse con una sola serie de precios. La ejecución de una orden, por ejemplo, depende del libro de órdenes o del pool de liquidez, de los costos, del tiempo de confirmación y de la estrategia de otros agentes. Por ello, el ecosistema CEX/DEX se aborda como un sistema complejo adaptativo, distribuido, no lineal y multi-agente.

\section{Planteamiento del problema}

Una proporción importante de los estudios de mercados cripto se concentra en predicción de precios o en backtesting simplificado. Estas aproximaciones pueden ser útiles como baselines, pero resultan insuficientes cuando omiten spread, profundidad, slippage, tarifas, gas, latencia, liquidaciones, fragmentación de liquidez, riesgo de oráculo y Maximum Extractable Value (MEV). Un modelo puede mostrar buen desempeño histórico bajo supuestos de ejecución ideales y fallar cuando se enfrenta a un cambio de régimen o a un shock de liquidez.

Existe además una brecha de reproducibilidad. Aunque los datos on-chain ofrecen trazabilidad y granularidad poco habituales en otros mercados, los experimentos suelen mezclar fuentes heterogéneas, transformaciones no documentadas y supuestos de ejecución implícitos. El resultado es difícil de auditar, repetir o extender. Esta situación es especialmente relevante en América Latina y Colombia, donde la adopción tecnológica avanza junto con una regulación heterogénea y una oferta todavía limitada de infraestructura académica para experimentar con estos sistemas.

La tesis propone una infraestructura experimental que conecte cuatro capas: adquisición y trazabilidad de datos; representación de la microestructura CEX/DEX; simulación de agentes heterogéneos; y evaluación reproducible de escenarios. El alcance inicial privilegia un simulador controlado y datasets documentados, antes de incorporar componentes de alta complejidad como aprendizaje por refuerzo, streaming o ejecución en tiempo real.

\section{Pregunta de investigación}

\textbf{¿Cómo diseñar una infraestructura experimental multi-agente, reproducible y auditable para modelar, simular y analizar ecosistemas financieros descentralizados híbridos CEX/DEX como sistemas complejos adaptativos?}

Esta pregunta se desarrolla mediante las siguientes preguntas específicas:

\begin{enumerate}
  \item ¿Qué variables de mercado, microestructura, protocolo y ejecución son necesarias para representar un escenario CEX/DEX de manera trazable?
  \item ¿Cómo modelar agentes heterogéneos --traders, proveedores de liquidez, arbitrajistas, liquidadores y agentes estratégicos simplificados-- sin confundir la simulación con una recomendación de inversión?
  \item ¿Qué fenómenos emergentes pueden observarse bajo shocks controlados de precio, liquidez, gas, latencia o colateral?
  \item ¿Qué valor aportan los modelos de series temporales y las representaciones multi-resolución como módulos de detección de régimen, frente a baselines transparentes?
  \item ¿Cómo evaluar reproducibilidad, sensibilidad, desempeño computacional y coherencia de la simulación?
\end{enumerate}

\section{Objetivos}

\subsection{Objetivo general}

Diseñar e implementar una infraestructura experimental multi-agente para la simulación y el análisis de ecosistemas financieros descentralizados híbridos CEX/DEX como sistemas complejos adaptativos, con trazabilidad de datos, evaluación reproducible y validación por escenarios.

\subsection{Objetivos específicos}

\begin{enumerate}
  \item Caracterizar la microestructura híbrida CEX/DEX y los agentes, eventos y dependencias relevantes para una representación experimental mínima.
  \item Diseñar un modelo de datos reproducible para fuentes CEX, DEX y on-chain, con separación explícita entre datos crudos, procesados y derivados.
  \item Implementar un simulador basado en agentes que incluya, al menos, traders, proveedores de liquidez, arbitrajistas y liquidadores simplificados.
  \item Incorporar baselines estadísticos y de aprendizaje automático para analizar regímenes, volatilidad y señales multi-resolución, sin asumir capacidad predictiva fuera de muestra.
  \item Evaluar escenarios de deterioro de liquidez, arbitraje con latencia, aumento de costos de ejecución, cascadas de liquidación y shocks exógenos.
  \item Documentar una matriz de riesgos operativos, de integridad de mercado y regulatorios como contexto de diseño, no como asesoría legal, tributaria o financiera.
  \item Publicar notebooks, documentación y artefactos de ejecución que permitan reproducir los experimentos del manuscrito.
\end{enumerate}

\section{Alcance de la versión 0.1.0}

La versión \versionTesis{} es un borrador de diseño. Consolida el problema, el marco conceptual, la arquitectura propuesta y el protocolo de evaluación que guiará la implementación. No presenta resultados empíricos definitivos, no opera con dinero real y no afirma rentabilidad de ninguna estrategia. Los resultados, figuras y comparaciones finales se incorporarán solo después de contar con fuentes verificadas, código versionado y ejecuciones reproducibles.

El primer incremento técnico se limita a un activo o par representativo, una fuente CEX, una fuente DEX u on-chain y un conjunto pequeño de escenarios controlados. Esta reducción de alcance busca asegurar que los supuestos sean inspeccionables antes de ampliar la plataforma.
